\chapter{Ascites (Abdominal Fluid Accumulation)}

\section*{Definition}
Ascites is abnormal fluid accumulation in the abdominal cavity. Cirrhosis and portal hypertension are common causes, but heart, kidney, cancerous, infectious, pancreatic, and other diseases can also cause it. New ascites needs medical evaluation because treatment depends on the cause.

\section*{When to See a Doctor}
\begin{itemize}
\item Progressive abdominal distension and increased abdominal girth
\item Rapid weight gain over days to weeks
\item Abdominal discomfort, fullness, or pain
\item Shortness of breath due to diaphragmatic compression
\item Nausea and early satiety
\item Umbilical hernia (from increased intra-abdominal pressure)
\item Lower extremity edema
\item Fever, new or worsening abdominal pain or tenderness, confusion, low blood pressure, or reduced urine (possible infection or organ complication)
\end{itemize}

\section*{How It Develops}
In cirrhosis, portal hypertension and changes in circulation cause the kidneys to retain sodium and water, allowing fluid to collect in the abdomen. Low albumin may contribute but is not the sole mechanism. Other causes involve heart or kidney disease, inflammation, cancer, lymphatic obstruction, or leakage of pancreatic or intestinal fluid.

\section*{Causes}
\begin{itemize}
\item Liver cirrhosis (most common cause)
\item Portal hypertension (non-cirrhotic causes such as portal vein thrombosis)
\item Congestive heart failure (right-sided heart failure)
\item Nephrotic syndrome
\item Peritoneal carcinomatosis (malignant ascites)
\item Pancreatitis (pancreatic ascites)
\item Tuberculous peritonitis
\item Budd-Chiari syndrome (hepatic vein thrombosis)
\item Myxedema (severe hypothyroidism)
\item Serositis due to connective tissue diseases (lupus)
\end{itemize}

\section*{Related Symptoms}
\begin{itemize}
\item Abdominal distension
\item Edema (peripheral swelling)
\item Shortness of breath
\item Jaundice
\item Abdominal pain
\item Nausea and vomiting
\item Umbilical hernia
\item Weight gain
\item Hepatomegaly or splenomegaly
\end{itemize}

\section*{Medical Evaluation}
\begin{itemize}
\item Physical examination and abdominal ultrasound, often with Doppler, to confirm fluid and assess the liver and portal or hepatic veins
\item Diagnostic paracentesis for all accessible new-onset ascites and for hospitalized people with cirrhosis and ascites, or sooner if fever, pain, gastrointestinal symptoms, kidney or liver deterioration, shock, or confusion occurs
\item Ascitic-fluid cell count, culture, protein, and albumin, together with serum albumin to calculate the serum-ascites albumin gradient (SAAG)
\item Liver function tests, coagulation profile, serum albumin, renal function
\item CT scan or MRI if malignancy is suspected
\end{itemize}

\section*{Treatment Options}
\begin{itemize}
\item Sodium restriction (less than 2 g/day)
\item Diuretic therapy (spironolactone with or without furosemide)
\item Large-volume paracentesis for tense or very symptomatic ascites; intravenous albumin is generally given when more than 5 L is removed to reduce circulatory complications
\item Transjugular intrahepatic portosystemic shunt (TIPS) for refractory ascites
\item Liver transplantation for end-stage liver disease
\item Antibiotics for spontaneous bacterial peritonitis
\end{itemize}

\section*{Self-Care and Home Management}
\begin{itemize}
\item Restrict dietary sodium to about 2 g per day if cirrhosis-related fluid retention is diagnosed, while avoiding malnutrition; fluid restriction is not routine unless clinically significant hyponatremia is present
\item Monitor daily weight, symptoms, and abdominal girth if the treating team recommends it, and report rapid changes
\item Elevate the head of the bed for comfort
\item Avoid alcohol completely if liver disease is present
\item Take prescribed diuretics as directed
\item Avoid NSAIDs which worsen sodium retention
\item Seek medical attention if fever, worsening pain, or confusion develops
\end{itemize}

\section*{References}
\begin{thebibliography}{99}
\bibitem{ascites:ref1} Moore KP, Aithal GP. ``Guidelines on the management of ascites in cirrhosis.'' \textit{Gut}. 2006;55(Suppl 6):vi1-vi12.
\bibitem{ascites:ref2} Runyon BA. ``Management of adult patients with ascites due to cirrhosis: update 2012.'' \textit{Hepatology}. 2013;57(4):1651-1653.
\bibitem{ascites:ref3} Gines P, Quintero E, et al. ``Compensated cirrhosis: natural history and prognosis.'' \textit{Semin Liver Dis}. 2015;35(3):278-289.
\bibitem{ascites:ref4} Schiff ER, Maddrey WC, et al. \textit{Schiff\'s Diseases of the Liver}, 12th ed. Wiley-Blackwell, 2017.
\bibitem{ascites:ref5} European Association for the Study of the Liver. ``EASL clinical practice guidelines on the management of ascites.'' \textit{J Hepatol}. 2018;69(2):406-460.
\bibitem{ascites:ref6} Biggins SW, Angeli P, Garcia-Tsao G, et al. Diagnosis, evaluation, and management of ascites, spontaneous bacterial peritonitis and hepatorenal syndrome: 2021 practice guidance. \textit{Hepatology}. 2021;74(2):1014--1048.
\end{thebibliography}
